\documentclass[12pt]{article}
\usepackage[utf8]{inputenc}
\usepackage[T1]{fontenc}
\usepackage{hyperref}
\title{Essono Language Resources for K-12 Curriculum Development}
\author{Compiled by Claude Sonnet 5.5 (claude-sonnet-5-5) for Andr\'as Kornai}
\date{}
\begin{document}
\maketitle

\section{Introduction}
ISO 639-3 sso is catalogued in linguistic sources as Sissano (also Sisano, Sinano, Sinama, Aissano, Siau), an Oceanic Austronesian language of the West Aitape Rural LLG, Sandaun Province, Papua New Guinea. Only Joshua Project uses the name ``Essono''; the search results suggest it is the same language, but that rests on a single mission-oriented source. Sissano is distinct from Arop (Arop-Sissano), a separate language with about 1,500 speakers. The 1998 tsunami nearly destroyed the community, and the language is classed as critically endangered. Speaker numbers fell from 4,780 (1990) to at most a few hundred.

\section{Primary Educational Resources}
No primers, textbooks or school programmes for Sissano/Essono were found.

\section{Linguistic and Reference Materials}
\begin{itemize}
\item Overview of the language: \url{https://en.wikipedia.org/wiki/Sissano_language}
\item Endangered Languages Project entry: \url{https://endangeredlanguages.com/lang/sso?hl=de}
\item Verbix endangered-language page: \url{https://docs.verbix.com/EndangeredLanguages/Sissano}
\end{itemize}
No grammar or dictionary turned up in these searches.

\section{Religious and Community Resources}
Joshua Project reports that no New Testament exists in the language and that only selected Bible portions have been translated: \url{https://joshuaproject.net/people_groups/14954/PP}. No community organisation was found.

\section{Digital and Audio Resources}
None identified in the search results.

\section{Curriculum Development Context}
Documentation is minimal and the community is small and disrupted by the tsunami. Any curriculum would depend on new documentation work with remaining speakers. Papua New Guinea's vernacular-education tradition is a possible model, but I found no source on it here. The language identity and name (Sissano or Essono) should be confirmed with the community.

\end{document}